% Options for packages loaded elsewhere
\PassOptionsToPackage{unicode}{hyperref}
\PassOptionsToPackage{hyphens}{url}
\PassOptionsToPackage{space}{xeCJK}
\documentclass[
  a4paper,
  11pt,
]{ctexart}
\usepackage{xcolor}
\usepackage[margin=2.35cm]{geometry}
\usepackage{amsmath,amssymb}
\setcounter{secnumdepth}{-\maxdimen} % remove section numbering
\setcounter{tocdepth}{2}
\usepackage{iftex}
\ifPDFTeX
  \usepackage[T1]{fontenc}
  \usepackage[utf8]{inputenc}
  \usepackage{textcomp} % provide euro and other symbols
\else % if luatex or xetex
  \usepackage{unicode-math} % this also loads fontspec
  \defaultfontfeatures{Scale=MatchLowercase}
  \defaultfontfeatures[\rmfamily]{Ligatures=TeX,Scale=1}
\fi
\usepackage{lmodern}
\ifPDFTeX\else
  % xetex/luatex font selection
  \ifXeTeX
    \usepackage{xeCJK}
    \setCJKmainfont[]{Microsoft YaHei}
  \fi
  \ifLuaTeX
    \usepackage[]{luatexja-fontspec}
    \setmainjfont[]{Microsoft YaHei}
  \fi
\fi
% Use upquote if available, for straight quotes in verbatim environments
\IfFileExists{upquote.sty}{\usepackage{upquote}}{}
\IfFileExists{microtype.sty}{% use microtype if available
  \usepackage[]{microtype}
  \UseMicrotypeSet[protrusion]{basicmath} % disable protrusion for tt fonts
}{}
\usepackage{setspace}
\makeatletter
\@ifundefined{KOMAClassName}{% if non-KOMA class
  \IfFileExists{parskip.sty}{%
    \usepackage{parskip}
  }{% else
    \setlength{\parindent}{0pt}
    \setlength{\parskip}{6pt plus 2pt minus 1pt}}
}{% if KOMA class
  \KOMAoptions{parskip=half}}
\makeatother
\usepackage{color}
\usepackage{fancyvrb}
\newcommand{\VerbBar}{|}
\newcommand{\VERB}{\Verb[commandchars=\\\{\}]}
\DefineVerbatimEnvironment{Highlighting}{Verbatim}{commandchars=\\\{\}}
% Add ',fontsize=\small' for more characters per line
\newenvironment{Shaded}{}{}
\newcommand{\AlertTok}[1]{\textcolor[rgb]{1.00,0.00,0.00}{\textbf{#1}}}
\newcommand{\AnnotationTok}[1]{\textcolor[rgb]{0.38,0.63,0.69}{\textbf{\textit{#1}}}}
\newcommand{\AttributeTok}[1]{\textcolor[rgb]{0.49,0.56,0.16}{#1}}
\newcommand{\BaseNTok}[1]{\textcolor[rgb]{0.25,0.63,0.44}{#1}}
\newcommand{\BuiltInTok}[1]{\textcolor[rgb]{0.00,0.50,0.00}{#1}}
\newcommand{\CharTok}[1]{\textcolor[rgb]{0.25,0.44,0.63}{#1}}
\newcommand{\CommentTok}[1]{\textcolor[rgb]{0.38,0.63,0.69}{\textit{#1}}}
\newcommand{\CommentVarTok}[1]{\textcolor[rgb]{0.38,0.63,0.69}{\textbf{\textit{#1}}}}
\newcommand{\ConstantTok}[1]{\textcolor[rgb]{0.53,0.00,0.00}{#1}}
\newcommand{\ControlFlowTok}[1]{\textcolor[rgb]{0.00,0.44,0.13}{\textbf{#1}}}
\newcommand{\DataTypeTok}[1]{\textcolor[rgb]{0.56,0.13,0.00}{#1}}
\newcommand{\DecValTok}[1]{\textcolor[rgb]{0.25,0.63,0.44}{#1}}
\newcommand{\DocumentationTok}[1]{\textcolor[rgb]{0.73,0.13,0.13}{\textit{#1}}}
\newcommand{\ErrorTok}[1]{\textcolor[rgb]{1.00,0.00,0.00}{\textbf{#1}}}
\newcommand{\ExtensionTok}[1]{#1}
\newcommand{\FloatTok}[1]{\textcolor[rgb]{0.25,0.63,0.44}{#1}}
\newcommand{\FunctionTok}[1]{\textcolor[rgb]{0.02,0.16,0.49}{#1}}
\newcommand{\ImportTok}[1]{\textcolor[rgb]{0.00,0.50,0.00}{\textbf{#1}}}
\newcommand{\InformationTok}[1]{\textcolor[rgb]{0.38,0.63,0.69}{\textbf{\textit{#1}}}}
\newcommand{\KeywordTok}[1]{\textcolor[rgb]{0.00,0.44,0.13}{\textbf{#1}}}
\newcommand{\NormalTok}[1]{#1}
\newcommand{\OperatorTok}[1]{\textcolor[rgb]{0.40,0.40,0.40}{#1}}
\newcommand{\OtherTok}[1]{\textcolor[rgb]{0.00,0.44,0.13}{#1}}
\newcommand{\PreprocessorTok}[1]{\textcolor[rgb]{0.74,0.48,0.00}{#1}}
\newcommand{\RegionMarkerTok}[1]{#1}
\newcommand{\SpecialCharTok}[1]{\textcolor[rgb]{0.25,0.44,0.63}{#1}}
\newcommand{\SpecialStringTok}[1]{\textcolor[rgb]{0.73,0.40,0.53}{#1}}
\newcommand{\StringTok}[1]{\textcolor[rgb]{0.25,0.44,0.63}{#1}}
\newcommand{\VariableTok}[1]{\textcolor[rgb]{0.10,0.09,0.49}{#1}}
\newcommand{\VerbatimStringTok}[1]{\textcolor[rgb]{0.25,0.44,0.63}{#1}}
\newcommand{\WarningTok}[1]{\textcolor[rgb]{0.38,0.63,0.69}{\textbf{\textit{#1}}}}
\usepackage{longtable,booktabs,array}
\usepackage{fancyhdr}
\usepackage{lastpage}
\setlength{\tabcolsep}{4pt}
\renewcommand{\arraystretch}{1.15}
\setlength{\LTleft}{0pt}
\setlength{\LTright}{0pt}
\pagestyle{fancy}
\fancyhf{}
\fancyhead[L]{MERBench MER2023 复现实验报告}
\fancyhead[R]{情感计算课程作业}
\fancyfoot[C]{\thepage\ / \pageref{LastPage}}
\setlength{\headheight}{14pt}
\usepackage{calc} % for calculating minipage widths
% Correct order of tables after \paragraph or \subparagraph
\usepackage{etoolbox}
\makeatletter
\patchcmd\longtable{\par}{\if@noskipsec\mbox{}\fi\par}{}{}
\makeatother
% Allow footnotes in longtable head/foot
\IfFileExists{footnotehyper.sty}{\usepackage{footnotehyper}}{\usepackage{footnote}}
\makesavenoteenv{longtable}
\setlength{\emergencystretch}{3em} % prevent overfull lines
\providecommand{\tightlist}{%
  \setlength{\itemsep}{0pt}\setlength{\parskip}{0pt}}
\usepackage{bookmark}
\IfFileExists{xurl.sty}{\usepackage{xurl}}{} % add URL line breaks if available
\urlstyle{same}
\hypersetup{
  pdftitle={MERBench MER2023 多模态情感识别复现实验报告},
  pdfauthor={情感计算课程作业},
  hidelinks,
  pdfcreator={LaTeX via pandoc}}

\title{MERBench MER2023 多模态情感识别复现实验报告}
\author{情感计算课程作业}
\date{2026 年 10 月}

\begin{document}
\maketitle
{\footnotesize
\tableofcontents
}
\newpage

\setstretch{1.18}
课程：情感计算\\
姓名：\_\_\_\_\_\_\_\_\_\_ 学号：\_\_\_\_\_\_\_\_\_\_\\
选题：MERBench: A Unified Evaluation Benchmark for Multimodal Emotion
Recognition\\
作者：Zheng Lian 等\\
实验日期：2026 年 10 月

\section{1
研究背景和问题定义}\label{1-ux7814ux7a76ux80ccux666fux548cux95eeux9898ux5b9aux4e49}

多模态情感识别需要同时利用语音、文本和视觉信息，预测样本的离散情绪类别以及连续情绪维度。不同研究往往使用不同的数据划分、特征提取器、融合模型和评价指标，因此即使报告了相似名称的结果，也不一定能够直接比较。MERBench
的目标是在统一数据、特征、训练和评价流程下比较多种多模态情感识别方法，并分析特征选择、融合策略、跨语料库迁移、噪声鲁棒性和预训练等因素。

本文定义的主要任务为：给定一个中文视频片段的音频特征 \(x^a\)、文本特征
\(x^t\) 和视觉特征 \(x^v\)，预测六分类情绪标签 \(y^e\) 以及连续 valence
标签 \(y^v\)。本次作业的主测试集为 MER-MULTI，分类标签为
neutral、angry、happy、sad、worried、surprise。

本复现严格锁定 MERBench，不使用 FER2013、FER+ 或其他数据集替代
MER-MULTI，也没有更换成其他论文。

\section{2
论文实验范围和本次复现范围}\label{2-ux8bbaux6587ux5b9eux9a8cux8303ux56f4ux548cux672cux6b21ux590dux73b0ux8303ux56f4}

论文完整的 MERBench
实验族包括单模态特征比较、多模态融合、跨语料库实验、标点鲁棒性、加性噪声鲁棒性、是否需要预训练，以及是否需要下游微调等内容。本次实验优先完成老师要求的
MER2023/MER-MULTI 核心部分：单模态基线、融合基线、结构改进和模态消融。

\subsection{2.1 MER-MULTI
论文结果行统计}\label{21-mer-multi-ux8bbaux6587ux7ed3ux679cux884cux7edfux8ba1}

以下统计只针对本项目整理的 MER-MULTI
参考结果文件，不代表论文全部实验族的总数量。

\begin{longtable}[]{@{}llrr@{}}
\toprule\noalign{}
论文表格 & 内容 & 结果行数 & 本地直接复现 \\
\midrule\noalign{}
\endhead
\bottomrule\noalign{}
\endlastfoot
Table II & high/medium 单模态 & 6 & 3 \\
Table III & 两种官方多模态特征组合 & 2 & 1 \\
Table IX & 五种融合模型的 high/medium 结果 & 10 & 5 \\
Table X & Top-k 特征选择 & 6 & 0 \\
\textbf{合计} & 论文参考表中的 MER-MULTI 结果 & \textbf{24} &
\textbf{9} \\
\end{longtable}

因此，本地严格对应论文结果行的覆盖率为 \(9/24=37.5\%\)。已经完成的 9
行包括 high 特征的三个单模态结果、一个 high 多模态特征组合结果，以及
Attention、LMF、MISA、TFN、MulT 五个 high 融合结果。Table IX 的 high
融合结果已经完成 5/5。

\subsection{2.2
本地实验配置统计}\label{22-ux672cux5730ux5b9eux9a8cux914dux7f6eux7edfux8ba1}

本地结果总表中目前有 13 个命名实验配置：5 个论文融合基线、3
个单模态实验、2 个改进实验和 3 个模态缺失消融实验。Attention
另使用固定超参运行了 seed 2026、2027、2028
三次统计，因此实际保存了额外的多 seed
运行结果。每个配置还分别记录交叉验证、MER-MULTI、MER-NOISE 和 MER-SEMI
的结果，这些测试划分输出不应被重复计数为新的独立模型。

\section{3 MERBench
方法和数据流程}\label{3-merbench-ux65b9ux6cd5ux548cux6570ux636eux6d41ux7a0b}

本地复现流程如下：

\begin{enumerate}
\def\labelenumi{\arabic{enumi}.}
\tightlist
\item
  使用官方 \texttt{label-6way.npz} 和转写文件读取 MER2023
  标签及数据划分。
\item
  使用作者提供的预提取音频、文本和视觉特征，避免在 8GB
  显存设备上重新提取大型模型特征。
\item
  将三种模态特征输入官方融合模型或本地结构改进模型。
\item
  同时优化离散情绪分类和 valence 回归任务。
\item
  在 Train\&Val 上进行五折交叉验证，并在 MER-MULTI、MER-NOISE、MER-SEMI
  上测试。
\item
  保存每次运行的配置、训练日志、每折最佳模型、预测结果、分类报告和混淆矩阵。
\end{enumerate}

MER2023 的本地划分如下。该划分与论文数据统计一致：Train\&Val 有 3373
个样本，MER-MULTI 有 411 个样本，MER-NOISE 有 412 个样本，MER-SEMI 有
834 个带离散标签样本；MER-SEMI 还包含大量未标注样本。

\begin{longtable}[]{@{}llrl@{}}
\toprule\noalign{}
split & 含义 & 样本数 & 本地使用方式 \\
\midrule\noalign{}
\endhead
\bottomrule\noalign{}
\endlastfoot
train\_corpus & Train\&Val & 3373 & 五折训练和验证 \\
test1\_corpus & MER-MULTI & 411 & 主要对照测试集 \\
test2\_corpus & MER-NOISE & 412 & 噪声鲁棒性测试输出 \\
test3\_corpus & MER-SEMI & 834 & 半监督相关测试输出 \\
\end{longtable}

MER-MULTI 的六类标签分布为：neutral 88、angry 105、happy 93、sad
50、worried 57、surprise 18。该分布不均衡，尤其 \texttt{surprise}
样本较少，因此报告同时使用 WAF 和 UAR 观察分类性能。

\section{4 特征和模型}\label{4-ux7279ux81f8ux548cux6a21ux578b}

\subsection{4.1
官方预提取特征}\label{41-ux5b98ux65b9ux9884ux63d0ux53d6ux7279ux5f81}

本地可验证读取的 high UTT 特征为：

\begin{longtable}[]{@{}lllr@{}}
\toprule\noalign{}
模态 & 特征 & 本地目录 & 单样本维度 \\
\midrule\noalign{}
\endhead
\bottomrule\noalign{}
\endlastfoot
audio & chinese HuBERT large & \texttt{chinese-hubert-large-UTT} &
1024 \\
text & Baichuan-13B-Base & \texttt{Baichuan-13B-Base-UTT} & 5120 \\
visual & CLIP ViT-L/14 & \texttt{clip-vit-large-patch14-UTT} & 768 \\
\end{longtable}

MulT 正式实验使用对应 FRA
特征：\texttt{chinese-hubert-large-FRA}、\texttt{Baichuan-13B-Base-FRA}
和 \texttt{clip-vit-large-patch14-FRA}。UTT 表示句子级特征，FRA
表示用于帧级对齐序列融合的特征。

\subsection{4.2 官方 Attention
融合基线}\label{42-ux5b98ux65b9-attention-ux878dux5408ux57faux7ebf}

对于每种模态，官方 Attention 模型先使用独立编码器得到统一隐空间表示：

\[h_a=E_a(x^a),\quad h_t=E_t(x^t),\quad h_v=E_v(x^v).\]

UTT 特征使用 MLP 编码器，FRA 特征使用 LSTM
编码器。三个模态表示拼接后输入注意力网络，得到三个模态的融合分数：

\[s=W_2\phi(W_1[h_a;h_t;h_v]),\]

然后按官方实现将三个隐向量进行加权融合，并通过两个输出头分别预测离散情绪和
valence：

\[\hat y^e=F_e(z),\quad \hat y^v=F_v(z).\]

训练目标为分类损失和回归损失之和，官方模型选择指标为：

\[M=WAF-0.25\times MSE.\]

\subsection{4.3 本地门控残差 Attention
改进}\label{43-ux672cux5730ux95e8ux63a7ux6b8bux5dee-attention-ux6539ux8fdb}

为了实现模型结构层面的改进，本文在官方 Attention
的三路编码器和双任务输出头基础上增加门控融合模块。设原始注意力分数为
\(s\)，首先计算归一化注意力权重：

\[\alpha=softmax(s).\]

再由拼接特征估计每个样本的模态可靠性：

\[r=\sigma(G([h_a;h_t;h_v])).\]

将两者相乘后重新归一化：

\[w_m=\frac{\alpha_m r_m}{\sum_j\alpha_jr_j+\epsilon}.\]

最终融合表示为：

\[z=\sum_{m\in\{a,t,v\}}w_mh_m+\beta R([h_a;h_t;h_v]),\]

其中 \(R\) 是轻量 MLP 残差支路，\(\beta\) 为可学习缩放参数，初始化为
0.1。该设计有两个目的：一是根据样本内容动态降低不可靠模态的影响，二是保留简单加权和不能表达的跨模态交互。该改进只增加融合层参数，不重新训练特征提取器，适合当前显存条件。

\subsection{4.4
类别加权改进}\label{44-ux7c7bux522bux52a0ux6743ux6539ux8fdb}

对于类别加权版本，分类损失使用基于训练集类别频数反比得到的 balanced
weight：

\[w_c=\frac{N}{C\cdot n_c},\]

其中 \(N\) 为训练样本数，\(C\) 为类别数，\(n_c\) 为第 \(c\)
类样本数。回归损失仍使用普通
MSE。该实验用于检验类别不均衡修正的效果，与门控残差结构改进分开运行。

\section{5
实验环境和运行设置}\label{5-ux5b9eux9a8cux73afux5883ux548cux8fd0ux884cux8bbeux7f6e}

\begin{longtable}[]{@{}ll@{}}
\toprule\noalign{}
项目 & 设置 \\
\midrule\noalign{}
\endhead
\bottomrule\noalign{}
\endlastfoot
操作系统 & Windows PowerShell \\
GPU & NVIDIA GeForce RTX 4060 Laptop GPU \\
显存 & 约 8GB \\
Python & 3.12.3 \\
PyTorch & 2.6.0+cu124 \\
CUDA & 12.4 \\
NumPy & 1.26.4 \\
scikit-learn & 1.7.2 \\
pandas & 2.2.2 \\
默认随机种子 & 2026 \\
UTT batch size & 32 \\
MulT FRA batch size & 16 \\
训练轮数 & 100 \\
交叉验证 & Train\&Val 五折 \\
\end{longtable}

本地环境与作者仓库中的旧版环境并不完全一致，因此本报告将结果称为``基于官方特征和官方代码结构的本地复现值''，不把它们写成论文原始报告值。

\section{6 评价指标}\label{6-ux8bc4ux4ef7ux6807ux6807}

\begin{itemize}
\tightlist
\item
  \textbf{WAF}：weighted F1，按类别样本数加权的 F1。
\item
  \textbf{ACC}：分类准确率。
\item
  \textbf{UAR}：macro recall，即各类别召回率的平均值，更关注少数类。
\item
  \textbf{MSE}：valence 回归均方误差，越低越好。
\item
  \textbf{模型选择指标}：\(WAF-0.25\times MSE\)。
\end{itemize}

MER-SEMI 的 valence 标签为 \texttt{-10}
占位值，因此本地结果中只报告其分类指标，不将其 MSE 与 MER-MULTI 或
MER-NOISE 的 MSE 比较。

\section{7
论文结果与本地复现结果}\label{7-ux8bbaux6587ux7ed3ux679cux4e0eux672cux5730ux590dux73b0ux7ed3ux679c}

下表主要列出论文 Table II、Table III 和 Table IX 中与当前 high
特征直接对应的结果。论文值来自本地保存的
\texttt{results/official\_mer2023\_mermulti.csv}，本地值来自 seed 2026
的真实运行结果。论文 WAF 和本地 WAF/ACC/UAR 均以百分数表示。

 {\scriptsize
\setlength{\tabcolsep}{2.5pt}
\begin{longtable}[]{@{}llrrrrrr@{}}
\toprule\noalign{}
实验 & 来源 & 论文 WAF & 论文 MSE & 本地 WAF & 本地 ACC & 本地 UAR & 本地 MSE \\
\midrule\noalign{}
\endhead
\bottomrule\noalign{}
\endlastfoot
Audio HUBERT-large & Table II & 69.91 & 0.97 & 68.59 & 68.61 & 63.56 &
0.9345 \\
Text Baichuan-13B & Table II & 62.54 & 1.63 & 60.68 & 60.83 & 55.36 &
1.6614 \\
Visual CLIP-large & Table II & 62.69 & 1.12 & 60.99 & 62.29 & 51.87 &
1.1297 \\
Attention high & Table IX / Table III & 83.03 & 0.62 & 80.45 & 80.78 &
75.78 & 0.6253 \\
LMF high & Table IX & 83.44 & - & 80.02 & 80.05 & 74.91 & 0.7023 \\
MISA high stable & Table IX & 82.42 & - & 80.00 & 80.05 & 75.11 &
0.6591 \\
TFN high & Table IX & 82.52 & - & 80.49 & 80.78 & 75.35 & 0.6934 \\
MulT high FRA & Table IX & 82.94 & - & 80.39 & 80.54 & 77.62 & 0.8183 \\
\end{longtable}
}\vspace{-0.4em}

\subsection{7.1 结果分析}\label{71-ux7ed3ux679cux5206ux6790}

\begin{enumerate}
\def\labelenumi{\arabic{enumi}.}
\tightlist
\item
  三个单模态本地结果与论文值处于同一数量级，audio 结果最接近论文 WAF。
\item
  多模态融合明显优于单模态，说明音频、文本和视觉之间存在互补信息。
\item
  五个 high 融合基线的本地 WAF 约为 80\%，比论文报告值低约 2.4 到 3.4
  个百分点。
\item
  Attention 的本地 MSE 0.6253 与论文 Table III high 特征组合的 0.62
  接近，但分类 WAF 低于论文值。
\item
  可能的差异来源包括论文多次运行均值和标准差、本地使用单
  seed、官方随机超参搜索没有完全复刻、MISA 使用了稳定超参，以及
  PyTorch/CUDA 版本不同。
\end{enumerate}

\subsection{7.2 Attention
多随机种子稳定性}\label{72-attention-ux591aux968fux673aux79cdux5b50ux7a33ux5b9aux6027}

固定 \texttt{configs/attention\_stable.yaml}，在 MER-MULTI 上运行三个
seed：

\begin{longtable}[]{@{}rrrrr@{}}
\toprule\noalign{}
seed & WAF & ACC & UAR & MSE \\
\midrule\noalign{}
\endhead
\bottomrule\noalign{}
\endlastfoot
2026 & 80.45 & 80.78 & 75.78 & 0.6253 \\
2027 & 83.06 & 82.97 & 80.19 & 0.6311 \\
2028 & 82.71 & 82.73 & 78.66 & 0.6157 \\
\textbf{均值 ± 标准差} & \textbf{82.07 ± 1.15} & \textbf{82.16 ± 0.98} &
\textbf{78.21 ± 1.83} & \textbf{0.6240 ± 0.0064} \\
\end{longtable}

Attention 的 WAF 标准差为 1.15 个百分点，UAR 标准差为 1.83
个百分点，说明单次 seed 结果存在一定波动。报告中的主要基线对照仍使用
seed 2026，三 seed 均值作为稳定性补充，不替代论文报告值。

\section{8
单模态、多模态和消融实验}\label{8-ux5355ux6a21ux6001ux591aux6a21ux6001ux548cux6d88ux878dux5b9eux9a8c}

以下为 MER-MULTI/test1 的本地结果：

\begin{longtable}[]{@{}llrrrr@{}}
\toprule\noalign{}
类型 & 实验 & WAF & ACC & UAR & MSE \\
\midrule\noalign{}
\endhead
\bottomrule\noalign{}
\endlastfoot
单模态 & audio only & 68.59 & 68.61 & 63.56 & 0.9345 \\
单模态 & text only & 60.68 & 60.83 & 55.36 & 1.6614 \\
单模态 & visual only & 60.99 & 62.29 & 51.87 & 1.1297 \\
多模态 & Attention high & 80.45 & 80.78 & 75.78 & 0.6253 \\
多模态 & LMF high & 80.02 & 80.05 & 74.91 & 0.7023 \\
多模态 & MISA high stable & 80.00 & 80.05 & 75.11 & 0.6591 \\
多模态 & TFN high & 80.49 & 80.78 & 75.35 & 0.6934 \\
多模态 & MulT high FRA & 80.39 & 80.54 & 77.62 & 0.8183 \\
消融 & Drop audio & 74.54 & 74.45 & 70.15 & 0.9129 \\
消融 & Drop text & 76.74 & 76.89 & 72.31 & 0.6763 \\
消融 & Drop visual & 75.22 & 75.18 & 71.08 & 0.8886 \\
\end{longtable}

Drop audio 后 WAF 比完整 Attention 下降 5.91 个百分点，Drop visual
后下降 5.23 个百分点，Drop text 后下降 3.71
个百分点。在当前特征和模型设置下，音频和视觉对离散分类的贡献更明显。Drop
text 的 MSE 仍为 0.6763，说明 valence 回归可能仍较依赖音频和视觉信息。

\section{9
算法改进和改进结果}\label{9-ux7b97ux6cd5ux6539ux8fdbux548cux6539ux8fdbux7ed3ux679c}

\subsection{9.1
类别加权改进}\label{91-ux7c7bux522bux52a0ux6743ux6539ux8fdb}

类别加权版本的 MER-MULTI 结果如下：

\begin{longtable}[]{@{}lrrrr@{}}
\toprule\noalign{}
实验 & WAF & ACC & UAR & MSE \\
\midrule\noalign{}
\endhead
\bottomrule\noalign{}
\endlastfoot
Attention high & 80.45 & 80.78 & 75.78 & 0.6253 \\
Attention + class weight & 81.56 & 81.51 & 78.45 & 0.6327 \\
\end{longtable}

相对于原始 Attention，WAF 提升 1.11 个百分点，ACC 提升 0.73
个百分点，UAR 提升 2.67 个百分点，MSE 增加
0.0074。该结果说明类别加权对于不均衡六分类任务有效，尤其改善宏平均召回；但分类权重调整不保证
valence 回归误差同时下降。

\subsection{9.2 门控残差
Attention}\label{92-ux95e8ux63a7ux6b8bux5dee-attention}

\begin{longtable}[]{@{}lrrrr@{}}
\toprule\noalign{}
实验 & WAF & ACC & UAR & MSE \\
\midrule\noalign{}
\endhead
\bottomrule\noalign{}
\endlastfoot
Attention high & 80.45 & 80.78 & 75.78 & 0.6253 \\
Gated residual Attention & 81.21 & 81.27 & 77.74 & 0.6369 \\
\end{longtable}

相对于原始 Attention，门控残差结构的 WAF 提升 0.76 个百分点，ACC 提升
0.49 个百分点，UAR 提升 1.96 个百分点，MSE 增加 0.0116。它优于原始
Attention，但 WAF 和 UAR
略低于类别加权版本。因此，本实验支持``门控残差结构有正向效果''，但不支持``门控残差结构在所有指标上优于类别加权''。

两种改进的侧重点不同：类别加权修改训练损失，针对类别不均衡；门控残差
Attention
修改融合结构，针对样本级模态可靠性和跨模态交互。为区分各自贡献，本轮没有直接叠加两种改进。

\section{10 误差分析}\label{10-ux8befux5deeux5206ux6790}

以 \texttt{Attention\ +\ class\ weight} 在 MER-MULTI 上的混淆矩阵为例：

\begin{longtable}[]{@{}lrrrrrr@{}}
\toprule\noalign{}
真实类别 \textbackslash{} 预测类别 & neutral & angry & happy & sad &
worried & surprise \\
\midrule\noalign{}
\endhead
\bottomrule\noalign{}
\endlastfoot
neutral & 68 & 3 & 7 & 4 & 6 & 0 \\
angry & 8 & 88 & 2 & 4 & 2 & 1 \\
happy & 8 & 0 & 84 & 0 & 0 & 1 \\
sad & 2 & 2 & 1 & 44 & 1 & 0 \\
worried & 10 & 0 & 0 & 7 & 40 & 0 \\
surprise & 1 & 1 & 2 & 1 & 2 & 11 \\
\end{longtable}

主要现象如下：

\begin{itemize}
\tightlist
\item
  \texttt{happy} 和 \texttt{angry}
  的识别相对稳定，主要因为测试样本数量较多且情绪线索较明显。
\item
  \texttt{surprise} 只有 18 个测试样本，类别加权后正确数为
  11，召回率约为 61.11\%，仍然容易受单个样本影响。
\item
  \texttt{worried} 容易被预测为 \texttt{neutral} 或
  \texttt{sad}，说明低唤醒负向情绪之间存在边界模糊。
\item
  类别加权后少数类表现有所改善，但整体 MSE 没有同步下降，体现了分类和
  valence 回归之间的任务权衡。
\end{itemize}

\section{11
复现产物和运行方法}\label{11-ux590dux73b0ux4ea7ux7269ux548cux8fd0ux884cux65b9ux6cd5}

官方代码和本地复现代码位于：

\begin{Shaded}
\begin{Highlighting}[]
\NormalTok{references/MERTools/MERBench}
\NormalTok{reproduction/MERBench}
\end{Highlighting}
\end{Shaded}

本地复现代码、配置、结果表和本报告源文件已公开在 GitHub：
\url{https://github.com/cai-jj21/merbench-mer2023-reproduction}

重要产物如下：

{\footnotesize
\begin{longtable}[]{@{}p{0.24\textwidth}p{0.66\textwidth}@{}}
\toprule\noalign{}
内容 & 路径 \\
\midrule\noalign{}
\endhead
\bottomrule\noalign{}
\endlastfoot
LaTeX 报告源文件 & \path{reports/merbench_report.tex} \\
报告 PDF & \path{reports/merbench_report.pdf} \\
GitHub 代码仓库 & \url{https://github.com/cai-jj21/merbench-mer2023-reproduction} \\
数据准备说明 & \path{DATA_PREP.md} \\
数据核验状态 & \path{DATA_STATUS.md} \\
运行 README & \path{README.md} \\
运行脚本 & \path{scripts/run_merbench_experiments.ps1} \\
官方结果参考表 & \path{results/official_mer2023_mermulti.csv} \\
本地结果总表 & \path{results/local_reproduction_summary.csv} \\
多 seed 汇总 & \path{results/attention_multiseed_summary.csv} \\
改进结果 & \path{results/gated_residual_attention_high_seed2026_formal/gated_residual_attention_high_seed2026_formal.csv} \\
改进模型代码 & \path{patches/toolkit/models/gated_attention.py} \\
日志和 checkpoint & \path{runs/} \\
\end{longtable}
}

从工作区根目录执行数据核验。为便于排版，命令中的脚本路径使用变量：

\begin{Verbatim}[fontsize=\scriptsize]
$run = ".\reproduction\MERBench\scripts\run_merbench_experiments.ps1"
$verify = ".\reproduction\MERBench\scripts\verify_mer2023_features.py"
python $verify `
  --data-root ".\reproduction\MERBench\data\mer2023-dataset-process" `
  --features chinese-hubert-large-UTT Baichuan-13B-Base-UTT `
  clip-vit-large-patch14-UTT
\end{Verbatim}

运行官方融合基线：

\begin{Verbatim}[fontsize=\scriptsize]
powershell -File $run -Stage attention
powershell -File $run -Stage lmf
powershell -File $run -Stage misa
powershell -File $run -Stage tfn
powershell -File $run -Stage mult_resume -Epochs 100 -BatchSize 16
\end{Verbatim}

运行改进、消融和汇总：

\begin{Verbatim}[fontsize=\scriptsize]
powershell -File $run -Stage improved -Epochs 100 -BatchSize 32
powershell -File $run -Stage classweight -Epochs 100 -BatchSize 32
powershell -File $run -Stage ablation -Epochs 100 -BatchSize 32
python .\reproduction\MERBench\scripts\build_summary_tables.py `
  --root .\reproduction\MERBench
\end{Verbatim}

\section{12
复现限制和未完成部分}\label{12-ux590dux73b0ux9650ux5236ux548cux672aux5b8cux6210ux90e8ux5206}

\begin{enumerate}
\def\labelenumi{\arabic{enumi}.}
\tightlist
\item
  当前 \texttt{feature.zip} 只包含 high UTT/FRA 特征目录，没有
  medium/low 特征池，因此无法完成论文 Table II、Table III、Table IX
  中对应的 medium 结果，也无法完成 Table X 的 Top-k 特征选择。
\item
  本地对 Attention 完成了三个 seed 的固定超参统计，LMF、MISA、TFN、MulT
  等其他正式基线主要使用 seed 2026，不能冒充论文的多次运行均值和标准差。
\item
  官方默认 MISA 随机超参在本机出现
  NaN，正式报告使用稳定超参版本，并保留默认运行失败日志。
\item
  本地没有重新提取 CLIP、HuBERT 和 Baichuan
  特征，因此复现范围是``基于作者官方预提取特征的融合训练与评测''。
\item
  本轮没有完成论文的跨语料库、标点鲁棒性、原始加性噪声生成流程、预训练和下游微调等完整实验族。MER-NOISE
  和 MER-SEMI
  的本地测试输出已经保存，但不能据此宣称完整复现论文全部鲁棒性和半监督实验。
\item
  本地环境为 Windows、Python 3.12、PyTorch 2.6，而作者 README
  中的环境较旧，版本差异可能影响数值结果。
\end{enumerate}

因此，本项目的准确表述是：\textbf{完成 MERBench/MER2023
核心融合实验的本地复现，以及一组结构改进和消融实验；尚未完成论文全部实验族的严格复现。}

\section{13 结论}\label{13-ux7ed3ux8bba}

本复现使用官方 MERBench 代码结构和官方 MER2023 预提取特征，在 RTX 4060
Laptop 8GB 环境下完成了可运行、可追溯的 MER-MULTI
核心实验。单模态实验表明，单独使用音频、文本或视觉的性能明显低于多模态融合；五种融合基线均表现出多模态互补优势。模态缺失消融显示，音频和视觉缺失对分类性能的影响大于文本缺失。

改进方面，类别加权版本取得了当前最高的 WAF 81.56 和 UAR 78.45；门控残差
Attention 取得 WAF 81.21 和 UAR 77.74，相比原始 Attention
有稳定的正向提升。两种改进分别从类别不均衡和模态可靠性两个角度改善模型，结果具有可解释性，但仍需要更多随机种子和完整特征池实验验证泛化性。

本地严格对应论文 MER-MULTI 参考结果完成 9/24
项，已满足本课程作业要求的论文基线、改进实验和消融实验要求；medium/low、Top-k、跨语料库及其他论文实验族明确列为尚未完成部分。

\section{14 参考文献}\label{14-ux53c2ux8003ux6587ux732e}

\begin{enumerate}
\def\labelenumi{\arabic{enumi}.}
\tightlist
\item
  Zheng Lian, Licai Sun, Yong Ren, Hao Gu, Haiyang Sun, Lan Chen, Bin
  Liu, Jianhua Tao. \emph{MERBench: A Unified Evaluation Benchmark for
  Multimodal Emotion Recognition}. arXiv:2401.03429, 2024.
  \url{https://arxiv.org/html/2401.03429}
\item
  MERBench
  官方代码仓库：\url{https://github.com/zeroQiaoba/MERTools/tree/master/MERBench}
\item
  MERBench 官方 README 中的 MER2023
  预提取特征入口：\url{https://pan.baidu.com/s/1l2yrWG3wXHjdRljAk32fPQ}；提取码：\texttt{uds2}
\end{enumerate}

\end{document}
